% !PRINTED 61
\weberpage{61}

Bezeichnen wir nun mit $\beta_1,\beta_2,\ldots,\beta_m$ ein zweites System von Indices, und bilden von der Function (7) die Ableitung $D_{\beta_1,\beta_2,\ldots,\beta_m}$, so ergiebt sich durch nochmalige Anwendung derselben Formeln (7) und (8):
\begin{equation}
 D_{\beta_1,\ldots,\beta_m}D_{\alpha_1,\ldots,\alpha_m}
 (x_1^{\mu_1}\cdots x_m^{\mu_m})
 =
 D_{\beta_1+\alpha_1,\ldots,\beta_m+\alpha_m}
 (x_1^{\mu_1}\cdots x_m^{\mu_m}),
\tag{9}
\end{equation}
und daraus folgt nach II. die allgemeine Gültigkeit des Satzes
\[
\text{III.}\qquad
D_{\beta_1,\ldots,\beta_m}D_{\alpha_1,\ldots,\alpha_m}\Phi
 =
D_{\beta_1+\alpha_1,\ldots,\beta_m+\alpha_m}\Phi,
\]
was eine Verallgemeinerung des Satzes (8), \S~13 ist, und man kann also die höheren Derivirten durch fortgesetzte Ableitung der niederen bilden.

Für die ersten Derivirten einer Function $\Phi$
\[
D_{1,0,\ldots,0}\Phi,\quad D_{0,1,\ldots,0}\Phi,\ldots,D_{0,0,\ldots,1}\Phi
\]
brauchen wir auch die kürzeren Zeichen
\[
\Phi'(x_1),\quad \Phi'(x_2),\ldots,\Phi'(x_m),
\]
ebenso für die zweiten
\[
D_{2,0,\ldots,0},\quad D_{1,1,\ldots,0},\ldots
\]
die Zeichen
\[
\Phi''(x_1,x_1),\quad \Phi''(x_1,x_2)=\Phi''(x_2,x_1),\ldots.
\]
Auch diese Bezeichnung lässt sich verallgemeinern und würde zu einem der Formel (4), \S~15 entsprechenden Ausdruck führen.

Wir erwähnen des häufigen Gebrauches wegen die Formeln für die quadratischen Formen besonders. Setzen wir
\begin{equation}
 \Phi(x)=\sum a_{i,k}x_i x_k,
\tag{10}
\end{equation}
worin $i,k$ von einander unabhängig die Reihe der Zahlen $1,2,\ldots,m$ durchlaufen und $a_{i,k}=a_{k,i}$ ist, so ist:
\begin{equation}
\begin{aligned}
 {1\over2}\Phi'(x_1)&=a_{1,1}x_1+a_{1,2}x_2+\cdots+a_{1,m}x_m,\\
 {1\over2}\Phi'(x_2)&=a_{2,1}x_1+a_{2,2}x_2+\cdots+a_{2,m}x_m,\\
 &\quad\cdots\\
 {1\over2}\Phi'(x_m)&=a_{m,1}x_1+a_{m,2}x_2+\cdots+a_{m,m}x_m,
\end{aligned}
\tag{11}
\end{equation}
und wir setzen noch
\begin{equation}
 \Phi(x,\xi)=\Phi(\xi,x)=\sum \xi_i\Phi'(x_i)=\sum x_i\Phi'(\xi_i).
\tag{12}
\end{equation}
Dann ist
\begin{equation}
 \Phi(x+\xi)=\Phi(x)+\Phi(x,\xi)+\Phi(\xi).
\tag{13}
\end{equation}
